\SourceLeaf{532}$-11713$, and the numbers represented by forms of the respective fundamental classes, for example $31\pdot\ 37\pdot\ 2$. As in my memoir in volume 14 of the Memoirs of the Royal Society of Sciences at Göttingen, fundamental classes here mean classes which, together with the classes in their periods, produce by composition every properly primitive class of the determinant once and only once.
\par\Needspace{5\baselineskip}The following entries have already been omitted:
\begin{longtable}{rcrl}
Hundred 9 & G. IV $\mathinner{.\,.\,.}$ & $3\pdot,$ & $-827[21::3]$ \\
26 & IV & $14$ & $-2587[24::11]$ \\
26 & VIII & $6$ & $-2564[56::3]$ \\
91 & I & $111$ & $-9059[117::5]$ \\
120 & IV & $32$ & $-11956{}^{*}2^{*}[36\pdot2::11\pdot49]$ \\
1 & I & $2$ & $+37[3::3]$ \\
2 & I & $2$ & $+101[3::4]$ \\
\end{longtable}
\par\Needspace{5\baselineskip}and added:
\begin{longtable}{rcrl}
Hundred 9 & G. IV $\mathinner{.\,.\,.}$ & $3\dotstwo$ & $-828[6\pdot2::31\pdot23]$ \\
26 & IV & $14$ & $-2586[28\pdot2::7\pdot2]$ \\
26 & VIII & $6$ & $-2565[12\pdot2\pdot2::7\pdot2\pdot5]$ \\
91 & I & $117$ & $-9059[117::5]$ \\
120 & IV & $32$ & $-11966{}^{*}2^{*}[32\pdot4::5\pdot83]$ \\
1 & I & $3$ & $+37[3::3]$ \\
2 & I & $3$ & $+101[3::4]$ \\
\end{longtable}
\par In the table for Hundred 3 and the last one on page 476, which in the manuscript occur together with a record of Hundreds 1 and 2 differing in appearance from the one printed here, the twelve divisions are designated, instead of I. Ordo unicus. 1; I. O. 2; I. O. 3; I. O. 4; II. Ordines duo. 1. 1; II. O. 1. 2; II. O. 2. 2; II. O. 3. 3; III. Ordines quatuor. 1. 1. 1. 1; III. O. 1. 1. 2. 2; III. O. 2. 2. 2. 2; IV. Ordines octo 1. 1. 1. 1. 1. 1. 1. 1, in the manner otherwise used here, as G. I. 1; G. I. 3; G. I. 5; G. I. 7; G. II. 1; G. II. 2; G. II. 3; G. II. 5; G. IV. 1; G. IV. 2; G. IV. 3; G. VIII. 1. Calculation gives, for example, 269. I. $3\ [3::4]$; 235. IV. $3\ [6\pdot2::3\pdot5]$; 401. I. $5\ [5::9]$; 577. I. $7\ [7::3]$; 727. II. $5\ [10::3]$. (Genera I in place of Genus I on page 469 is a printing error).
\par\Needspace{5\baselineskip}Because of printing errors, omit:
\begin{longtable}{rcrl}
Hundred 27. & G. IV $\dotsfour$ & $16\mathinner{.\,.\,.}$ & $-2624{}^{*}3^{*}$ \\
93 & IV & $16$ & $-9216$ \\
118 & VIII & $4$ & $-11713{}^{*}3^{*}$ \\
\end{longtable}
\par\Needspace{5\baselineskip}and add:
\begin{longtable}{rcrl}
Hundred 27. & G. IV $\dotsfour$ & $16\mathinner{.\,.\,.}$ & $-2624{}^{*}2^{*}[16\pdot4::3\pdot16]$ \\
93 & IV & $16$ & $-9216{}^{*}2^{*}[16\pdot4::5\pdot9]$ \\
118 & VIII & $4$ & $-11713{}^{*}2^{*}[4\pdot4\pdot2::31\pdot37\pdot2]$ \\
\end{longtable}
\par\Needspace{5\baselineskip}According to my calculation, also omit:
\begin{longtable}{rcrl}
Hundred 10. & G. II $\dotsfour$ & $9\mathinner{.\,.\,.}$ & $-972[6\pdot3::7\pdot13]$ \\
17 & IV & $4$ & $-1660[10\pdot2::11\pdot5]$ \\
20 & IV & $12$ & $-1982[24::3]$ \\
21 & IV & $6$ & $-2096[30\pdot2::3\pdot4]$ \\
23 & IV & $9$ & $-2221[18::10]$ \\
24 & IV & $12$ & $-2376[12\pdot2\pdot2::5\pdot8\pdot8]$ \\
29 & IV & $9$ & $-2887[25::8]$ \\
61 & IV & $7$ & $-6028[12\pdot2::13\pdot4]$ \\
96 & VIII & $13$ & $-9594[20\pdot2\pdot2::31\pdot2\pdot13]$ \\
118 & IV & $25$ & $-11780[16\pdot4\pdot2::3\pdot8\pdot19]$ \\
118 & VIII & $16$ & $-11780[16\pdot4\pdot2::3\pdot8\pdot19]$ \\
119 & VIII & $16$ & $-11840[24\pdot2\pdot2::5\pdot9\pdot7]$ \\
\noalign{\SourceLeaf{533}}
Thousand I & G. II $\mathinner{.\,.\,.}$ & $3\mathinner{.\,.\,.}$ & $-541[10::11]$ \\
I & II & $4$ & $-415[10::13]$ \\
I & II & $8$ & $-527[18::3]$ \\
I & II & $8$ & $-722[18::3]$ \\
I & II & $9$ & $-194[20::5]$ \\
I & II & $9$ & $-459[6\pdot3::5\pdot9]$ \\
I & II & $9$ & $-972[6\pdot3::7\pdot13]$ \\
I & II & $11$ & $-842[26::13]$ \\
I & IV & $3$ & $-784[8\pdot2::5\pdot4]$ \\
I & IV & $4$ & $-532[4\pdot2::13\pdot7]$ \\
I & IV & $5$ & $-425[12\pdot2::3\pdot17]$ \\
I & IV & $5$ & $-608[12\pdot2::13\pdot27]$ \\
I & IV & $5$ & $-629[18\pdot2::5\pdot2]$ \\
III & II & $15$ & $-2578[16::13]$ \\
X & I & $111$ & $-9059[117::5]$ \\
\multicolumn{2}{l}{forms $-(15n+13)$ IV} & $4$ & $-2788{}^{*}2^{*}[8\pdot2::19\pdot17]$ \\
\end{longtable}
\par\Needspace{5\baselineskip}and add:
\begin{longtable}{rcrl}
Hundred 10. & G. II $\dotsfour$ & $9$ & $-972{}^{*}3^{*}[6\pdot3::7\pdot13]$ \\
17 & IV & $12$ & $-1700[24\pdot2::3\pdot17]$ \\
20 & IV & $12$ & $-1937[24\pdot2::7\pdot2]$ \\
21 & IV & $6$ & $-2097[12\pdot2::47\pdot2]$ \\
23 & IV & $9$ & $-2224[18\pdot2::5\pdot16]$ \\
24 & IV & $12$ & $-2366[24\pdot2::3\pdot2]$ \\
29 & IV & $9$ & $-2885[18\pdot2::3\pdot5]$ \\
61 & IV & $6$ & $-6028[12\pdot2::13\pdot4]$ \\
96 & VIII & $13$ & $-9546[26\pdot2\pdot2::5\pdot3\pdot37]$ \\
118 & IV & $25$ & $-11750[50\pdot2::3\pdot47]$ \\
118 & VIII & $16$ & $-11780{}^{*}2^{*}[16\pdot4\pdot2::3\pdot8\pdot19]$ \\
119 & VIII & $16$ & $-11840{}^{*}2^{*}[16\pdot4\pdot2::3\pdot16\pdot5]$ \\
Thousand I. & G. II $\dotsfour$ & $5\mathinner{.\,.\,.}$ & $-415[10::13]$ \\
I & II & $5$ & $-541[10::11]$ \\
I & II & $9$ & $-459{}^{*}3^{*}[6\pdot3::5\pdot9]$ \\
I & II & $9$ & $-527[18::3]$ \\
I & II & $9$ & $-722[18::3]$ \\
I & II & $9$ & $-972{}^{*}3^{*}[6\pdot3::7\pdot13]$ \\
I & II & $10$ & $-194[20::5]$ \\
I & II & $13$ & $-842[26::13]$ \\
I & IV & $2$ & $-532[4\pdot2::13\pdot7]$ \\
I & IV & $4$ & $-784[8\pdot2::5\pdot4]$ \\
I & IV & $6$ & $-425[12\pdot2::3\pdot17]$ \\
I & IV & $6$ & $-608[12\pdot2::13\pdot27]$ \\
I & IV & $9$ & $-629[18\pdot2::5\pdot2]$ \\
III & II & $15$ & $-2518[30::19]$ \\
X & I & $117$ & $-9059[117::5]$ \\
\multicolumn{2}{l}{forms $-(15n+13)$ IV} & $4$ & $-2788[8\pdot2::19\pdot17]$ \\
\end{longtable}
\par The tables \textit{for Circle Division} give, for 2452 numbers of the forms $aa+1$, $aa+4$, $aa+9$, $\dotsfour aa+81$, all the odd prime divisors $p$ beside the corresponding $a$, in those cases where all the prime divisors are less than 200; only then are $aa+1$ and so on called factorable.
\par To make the tables easier to use, \textsc{Gauss} has constructed, for each table from which the complete factorizations of numbers of one of the particular forms can be determined, an auxiliary table giving beside each prime $p$ those numbers $a$ whose square, increased by 1 or $4\mathinner{.\,.\,.}$, has $p$ as its greatest prime divisor.
\par The chief purpose of the tables is to facilitate the accurate calculation of arcs whose cotangents are given rational numbers. With their help, arcs for small cotangents can first be combined from arcs for large cotangents, reducing to a very small amount the calculations still needed for the series expressing arcs in terms of their cotangents. The developments towards this end in the manuscript remains are not extensive; the following is the one carried furthest. In it
\[
\begin{array}{l}
[2]\ [5]\ [13]\ [17]\ [29]\ [37]\ [41]\ [53]\ [61]\mathinner{.\,.\,.}[197]\\
(18)\ (57)\ (239)\ \left(\frac{79}{3}\right)\mathinner{.\,.\,.}
\end{array}
\]
denote the arcs with cotangents
\[
1\quad2\quad\frac{3}{2}\quad4\quad\frac{5}{2}\quad6\quad\frac{5}{4}\quad\frac{7}{2}\quad\frac{6}{5}\mathinner{.\,.\,.}14\quad18\quad57\quad239\quad\frac{79}{3}\mathinner{.\,.\,.}
\]
Using the tables to factor $18+i$, $57+i$, $239+i$ into their complex prime factors gives
\[
\begin{aligned}
(18)&=2[2]-2[5]-[13]\\
(57)&=-[2]+3[5]-[13]\\
(239)&=3[2]\phantom{-2[5]}-4[13]
\end{aligned}
\]
